\documentclass[11pt]{article}
\usepackage{explorer}
\newcommand{\ExplorerLanguage}{EN}
\newcommand{\ExplorerNumber}{08}
\newcommand{\ExplorerKind}{Worked solutions}
\begin{document}
\ExplorerTitle{Upper high school · probability and sequences}{Waiting for the first success}
\ExplorerSection{Exercise 1}
With \(p=\frac{1}{4}\), find \(P(T=3)\).\par
\((\frac{3}{4})^{2}(\frac{1}{4})=\frac{9}{64}\).\par
\ExplorerSection{Exercise 2}
With \(p=\frac{1}{4}\), find \(P(T\leq 3)\).\par
\(1-P(T>3)=1-(\frac{3}{4})^{3}=\frac{37}{64}\).\par
\ExplorerSection{Exercise 3}
For \(0<p<1\), find \(P(T>m+n|T>m)\).\par
The ratio of tail probabilities is \((1-p\))ᵐ⁺ⁿ/\((1-p\))ᵐ=\((1-p)^{n}\). After m failures, the remaining waiting-time distribution is unchanged.\par
\ExplorerSection{Exercise 4}
Success has probability 0.2 at each independent trial. What is the mean waiting time? Must the fifth trial succeed?\par
\(E(T)=5\) trials. No: every trial still has success probability 0.2, independently of previous failures.\par
\ExplorerSection{Problem: complete reasoning}
Roll a fair die until the first 6, but stop after 10 rolls in any case. Let N be the number of rolls actually made. Find \(P(N=10)\), the probability of obtaining at least one 6, and E(N). For the mean, use \(E(N)=\sum  P(N>n)\), from \(n=0\) to 9.\par
\(N=10\) when the first nine rolls fail, regardless of the tenth: \(P(N=10)=(\frac{5}{6})^{9}\). The probability of at least one 6 is \(1-(\frac{5}{6})^{10}\). For \(0\leq n\leq 9, P(N>n)=(\frac{5}{6})^{n}\), giving \(E(N)=6[1-(\frac{5}{6})^{10}]\), about 5.031 rolls. Truncation reduces the expectation below 6.\par
\end{document}
